\section{消息历史的运行时管理}

\subsection{多轮交互的 Context 管理}

在预注入的 4 条消息之后，Codex 进入真正的用户-Agent 多轮交互阶段：

\begin{enumerate}
    \item 用户发出一个真实的 Query
    \item Agent Loop 开始执行多轮（multi-turn）交互
    \item 交互过程中包含工具调用（Function Call）及其返回结果
    \item 一个 Turn 结束
    \item 用户发出新的 Query，开始下一个 Turn
\end{enumerate}

\begin{importantbox}{Compact 指令的作用}
当消息历史积累过多时，可以在 Codex CLI 中执行 \texttt{compact} 指令。Compact 压缩的是多轮对话部分的内容，但会保留预注入的系统消息。这是 Context Engineering 中管理上下文窗口的关键手段。
\end{importantbox}

\subsection{Skills 的渐进式加载}

\begin{knowledgebox}{Skills 的两阶段加载机制}
Skills 采用渐进式加载（Progressive Loading）策略：
\begin{enumerate}
    \item \textbf{阶段一（注册）}：在预注入的 User 消息中，只暴露每个 Skill 的名称和简短描述
    \item \textbf{阶段二（加载）}：当远端模型决定使用某个具体 Skill 时，通过 Function Call 执行 \texttt{cat} 命令，将该 Skill 对应的 skills.md 文件完整内容作为 Function Output 加载到 Context 中
\end{enumerate}
这种设计避免了一次性加载所有 Skills 导致的 Context 浪费。
\end{knowledgebox}

\subsection{本章小结}

Codex 的消息历史在运行时持续增长，通过 compact 指令进行上下文压缩。Skills 采用渐进式加载，先暴露摘要，按需加载完整内容，是一种精巧的 Context 管理策略。


